\section{Tres objetivos de predicción equivalentes}

Dado que la media posterior real $\tilde{\mu}_t(x_t, x_0)$ es una función lineal de $x_t$ y $x_0$, y existe una relación determinista entre $x_0$, $x_t$ y el ruido $\epsilon$, podemos reparametrizar la predicción de la media en tres objetivos de predicción equivalentes.

\subsection{Enfoque 1: Predicción de la media (Mean Prediction)}

La forma más directa consiste en que la red neuronal prediga directamente la media de la transición inversa:

$$\mu_\theta(x_t, t) \approx \tilde{\mu}_t(x_t, x_0)$$

La red acepta $x_t$ y $t$ como entrada y produce un vector de medias con la misma dimensión que $x_t$.

La función de pérdida de entrenamiento es:

$$L_{\text{mean}} = \mathbb{E}_{t, x_0, x_t}\left[\left\| \tilde{\mu}_t(x_t, x_0) - \mu_\theta(x_t, t) \right\|^2\right]$$

\begin{knowledgebox}{Características de la predicción de la media}
La predicción de la media es la forma de parametrización más natural; corresponde directamente al objetivo de optimización derivado del ELBO. Sin embargo, la escala de la media $\tilde{\mu}_t$ varía con el paso temporal $t$ (ya que es una combinación ponderada de $x_t$ y $x_0$, donde los pesos dependen de $t$), lo que puede provocar un entrenamiento poco estable.
\end{knowledgebox}

\subsection{Enfoque 2: Predicción de la muestra limpia (Clean Sample Prediction / $x_0$-Prediction)}

Dado que $\tilde{\mu}_t(x_t, x_0)$ es una función lineal de $x_t$ y $x_0$, y $x_t$ se conoce durante la inferencia, podemos pedirle a la red que prediga $x_0$ en lugar de predecir directamente la media:

$$\hat{x}_0 = f_\theta(x_t, t)$$

A continuación, calculamos la media mediante la fórmula:

$$\mu_\theta(x_t, t) = \frac{\sqrt{\bar{\alpha}_{t-1}}\, \beta_t}{1 - \bar{\alpha}_t}\, f_\theta(x_t, t) + \frac{\sqrt{\alpha_t}\, (1 - \bar{\alpha}_{t-1})}{1 - \bar{\alpha}_t}\, x_t$$

La función de pérdida de entrenamiento equivalente (ignorando los términos de ponderación) es:

$$L_{x_0} = \mathbb{E}_{t, x_0, \epsilon}\left[\left\| x_0 - f_\theta(x_t, t) \right\|^2\right]$$

\begin{importantbox}{Perspectiva única de la predicción $x_0$}
La predicción $x_0$ ofrece una intuición de muestreo atractiva: en cada paso intermedio $t$ del proceso inverso, la red \textbf{predice el valor esperado de la salida final $x_0$}, y luego decide cómo muestrear el siguiente paso $x_{t-1}$ basándose en ello. Esto significa que el proceso de muestreo de los modelos de difusión es un proceso de \textbf{refinamiento iterativo} (iterative refinement): en cada paso se predice el resultado final y, a continuación, la predicción se corrige gradualmente en lugar de simplemente generar una salida directa.
\end{importantbox}

\begin{warningbox}{La predicción $x_0$ no es generación en un solo paso}
Aunque en cada paso se predice $x_0$, no utilizamos directamente esta predicción como salida final. En su lugar, calculamos la media de la transición inversa con $\hat{x}_0$ y luego muestreamos $x_{t-1}$; posteriormente, predecimos nuevamente $\hat{x}_0$ sobre $x_{t-1}$, iterando así hasta llegar a $x_0$. La predicción de $\hat{x}_0$ en cada paso se vuelve gradualmente más precisa.
\end{warningbox}

\subsection{Enfoque 3: Predicción del ruido (Noise Prediction / $\epsilon$-Prediction)}

Esta es la forma de parametrización adoptada en el artículo DDPM y también es la más utilizada en la práctica.

Recordemos la reparametrización de la distribución de salto hacia adelante:

$$x_t = \sqrt{\bar{\alpha}_t}\, x_0 + \sqrt{1 - \bar{\alpha}_t}\, \epsilon, \quad \epsilon \sim \mathcal{N}(0, I)$$

De aquí podemos expresar $x_0$ como una función de $x_t$ y $\epsilon$:

$$x_0 = \frac{x_t - \sqrt{1 - \bar{\alpha}_t}\, \epsilon}{\sqrt{\bar{\alpha}_t}}$$

Sustituyendo esto en la expresión de $\tilde{\mu}_t$, obtenemos una media expresada en términos de $x_t$ y $\epsilon$:

$$\tilde{\mu}_t(x_t, \epsilon) = \frac{1}{\sqrt{\alpha_t}}\left(x_t - \frac{\beta_t}{\sqrt{1 - \bar{\alpha}_t}}\, \epsilon\right)$$

Por lo tanto, podemos pedirle a la red que prediga el ruido $\epsilon$ en lugar de la media o $x_0$:

$$\hat{\epsilon} = \epsilon_\theta(x_t, t)$$

A continuación, calculamos la media mediante la fórmula:

$$\mu_\theta(x_t, t) = \frac{1}{\sqrt{\alpha_t}}\left(x_t - \frac{\beta_t}{\sqrt{1 - \bar{\alpha}_t}}\, \epsilon_\theta(x_t, t)\right)$$

\begin{importantbox}{Objetivo de entrenamiento simplificado de DDPM}
El hallazgo central de Ho et al. (2020): tras eliminar los términos de ponderación, el objetivo de entrenamiento de la predicción del ruido se reduce a una forma extremadamente simple:
$$L_{\text{simple}} = \mathbb{E}_{t, x_0, \epsilon}\left[\left\| \epsilon - \epsilon_\theta(x_t, t) \right\|^2\right]$$
donde $t \sim \text{Uniform}(\{1, \dots, T\})$, $x_0 \sim q(x_0)$, $\epsilon \sim \mathcal{N}(0, I)$ y $x_t = \sqrt{\bar{\alpha}_t}\, x_0 + \sqrt{1 - \bar{\alpha}_t}\, \epsilon$.

Esto es: \textbf{dada una imagen con ruido, predecir el ruido añadido, minimizando el error L2}.
\end{importantbox}

\subsection{Equivalencia y selección entre las tres parametrizaciones}

Las tres parametrizaciones son matemáticamente completamente equivalentes, ya que se pueden convertir unas a otras:

\begin{itemize}
    \item De $\hat{x}_0$ se puede calcular $\hat{\epsilon} = \frac{x_t - \sqrt{\bar{\alpha}_t}\, \hat{x}_0}{\sqrt{1-\bar{\alpha}_t}}$
    \item De $\hat{\epsilon}$ se puede calcular $\hat{x}_0 = \frac{x_t - \sqrt{1-\bar{\alpha}_t}\, \hat{\epsilon}}{\sqrt{\bar{\alpha}_t}}$
    \item De $\hat{x}_0$ o $\hat{\epsilon}$ se puede calcular $\hat{\mu}$
\end{itemize}

\begin{table}[H]
\centering
\caption{Comparación de los tres objetivos de predicción}
\begin{tabular}{llll}
\toprule
\textbf{Forma de parametrización} & \textbf{Objetivo de predicción} & \textbf{Función de pérdida} & \textbf{Observaciones} \\
\midrule
Predicción de la media & $\tilde{\mu}_t$ & $\|\tilde{\mu}_t - \mu_\theta\|^2$ & Más natural pero depende de la escala en $t$ \\
Predicción $x_0$ & $x_0$ & $\|x_0 - f_\theta(x_t,t)\|^2$ & Intuición de refinamiento iterativo \\
Predicción $\epsilon$ & $\epsilon$ & $\|\epsilon - \epsilon_\theta(x_t,t)\|^2$ & Predeterminado en DDPM, más estable \\
\bottomrule
\end{tabular}
\end{table}

\begin{importantbox}{¿Por qué la predicción del ruido es la más común?}
En la práctica, la predicción del ruido ($\epsilon$-prediction) es la forma de parametrización ampliamente utilizada; las razones principales son la \textbf{estabilidad del entrenamiento}:
\begin{enumerate}
    \item \textbf{Normalización del valor objetivo}: el ruido $\epsilon \sim \mathcal{N}(0,I)$ es una muestra de una distribución normal estándar, que posee naturalmente características de escala óptimas con media cero y varianza unitaria. En cambio, la escala de $x_0$ depende del conjunto de datos, y la escala de $\tilde{\mu}_t$ varía con $t$.
    \item \textbf{Consistencia en la escala del gradiente}: dado que el objetivo del ruido tiene las mismas características estadísticas en todos los pasos temporales, la escala de los gradientes es más consistente entre diferentes pasos temporales, lo cual favorece la estabilidad del entrenamiento.
    \item \textbf{Sin normalización adicional}: para los datos $x_0$, generalmente se requiere normalizarlos a un rango específico; mientras que el ruido $\epsilon$ ya está normalizado por sí mismo.
\end{enumerate}
\end{importantbox}

\begin{warningbox}{Diferentes parametrizaciones son adecuadas para diferentes escenarios}
Aunque la predicción del ruido ofrece los mejores resultados en la mayoría de los escenarios, la predicción $x_0$ puede ser más ventajosa en aplicaciones específicas, por ejemplo: edición de imágenes (necesita estimar directamente la imagen limpia), generación de video (necesita mantener la consistencia temporal), entre otras. La parametrización v-prediction posterior combina las ventajas de ambas. Al elegir la forma de parametrización, se deben considerar los requisitos específicos de la tarea.
\end{warningbox}

\subsection{Resumen de la sección}

Los tres objetivos de predicción son matemáticamente equivalentes, pero muestran comportamientos distintos durante el entrenamiento real. La predicción del ruido se ha convertido en la opción por defecto debido a sus buenas características de normalización y estabilidad del entrenamiento. La predicción de muestras limpias ofrece una perspectiva intuitiva de refinamiento iterativo. Las tres formas son intercambiables; en implementaciones prácticas, esta relación de conversión suele aprovecharse durante la inferencia.

%% ========== Sección 6: Algoritmo de entrenamiento ==========
